\documentclass[a4paper]{scrartcl}
\usepackage{amsmath,amssymb,amsthm}
\usepackage{bm}
\usepackage{biblatex}
\usepackage{caption}
\usepackage{float}
\usepackage[colorlinks=true, allcolors=black, urlcolor=cyan]{hyperref}
\usepackage{graphicx}
\usepackage[framed,numbered]{matlab-prettifier}
\usepackage{mathtools}
\usepackage{mcode}
\usepackage{physics}
\usepackage{siunitx}
\usepackage{cancel}

\renewcommand{\thesection}{\arabic{section}}
\renewcommand{\thesubsection}{\arabic{section}.\alph{subsection}}

\title{Assignment 3}
\subtitle{ELEC 442 - Introduction to Robotics}
\author{Sondre Myrberg (81113433) \and Ola Helbaek (68776772)}


\setlength{\parindent}{0pt} % Disable indentation
\mathtoolsset{showonlyrefs} % Only show equation numbers on referenced equations

\def\undertilde#1{\mathord{\vtop{\ialign{##\crcr
$\hfil\displaystyle{#1}\hfil$\crcr\noalign{\kern1.5pt\nointerlineskip}
$\hfil\widetilde{}\hfil$\crcr\noalign{\kern1.5pt}}}}} %Lousy way, but need this undertilde
\newcommand{\me}[1]{\mathrm{e}^{#1}}


\begin{document}

\hypersetup{pageanchor=false}
\begin{titlepage}
    \maketitle
    \vfill
    \vfill
    \vfill
    \vfill
    \includegraphics[width=0.95\textwidth]{../../ubc_logo.pdf}
    \vfill
    \vfill
\end{titlepage}
\hypersetup{pageanchor=true}

\section{Computation of representation}

As we have the relationship between the initial frame $\underline{C}_i$ and the final frame $\underline{C}_f$ given as
\begin{equation}
    \begin{aligned}
        \underline{C}_f &= \underline{\underline{R}}\underline{C}_i\\
        &= \me{\theta \underline{\bm{s}}\cross}\underline{C}_i
    \end{aligned}
\end{equation}
we could rearrange this as
\begin{equation}
    \begin{aligned}
        \underline{C}_f &= \underline{C}_i R\\
         &= \underline{C}_i \me{\theta\bm{s}\cross}
    \end{aligned}
\end{equation}
where
\begin{equation}
    R = \begin{bmatrix}
        r_{11} & r_{12} & r_{13} \\  r_{21} & r_{22} & r_{23} \\  r_{31} & r_{32} & r_{33}
    \end{bmatrix}
\end{equation}
as given in the text. This gives us the relationship
\begin{equation}
    \begin{aligned}
        \me{\theta\bm{s}\cross} = \begin{bmatrix}
        r_{11} & r_{12} & r_{13} \\  r_{21} & r_{22} & r_{23} \\  r_{31} & r_{32} & r_{33}
    \end{bmatrix}
    \end{aligned}
\end{equation}
which we want to solve for $\bm{s}$. Firstly use that
\begin{equation}
    \begin{aligned}
         \me{\theta\bm{s}\cross} &= \sum_{n=0}^{\infty}\frac{(\theta \bm{s}\cross)^n}{n!}\\
         &= \bm{I} + (\theta \bm{s}\cross) + \frac{1}{2!}(\theta \bm{s}\cross)^2 + \hdots
    \end{aligned}
\end{equation}
and with the property that $(\bm{s}\cross)^3 = - (\bm{s}\cross)$ for a skew-symmetric matrix we get
\begin{equation}
    \begin{aligned}
        \me{\theta\bm{s}\cross} &= \bm{I} + \underbrace{\left(\theta - \frac{\theta^3}{3!} + \frac{\theta^5}{5!} - \hdots \right)}_{= \sin\theta}(\bm{s}\cross)
        + \underbrace{\left(\frac{\theta^2}{2!} - \frac{\theta^4}{4!} +\frac{\theta^6}{6!} - \hdots \right)}_{= 1- \cos\theta}(\bm{s}\cross)^2\\
        &= \bm{I} + \sin\theta(\bm{s}\cross) + (1-\cos\theta)(\bm{s}\cross)^2
        = R
    \end{aligned}
\end{equation}
With
\begin{equation}
    \bm{s}\cross = \begin{bmatrix}
        0 & -s_3 & s_2\\
        s_3 & 0 & -s_1\\
        -s_2 & s_1 & 0
    \end{bmatrix}
\end{equation}
and
\begin{equation}
    (\bm{s}\cross)^2 = \left[\begin{array}{ccc}  - {s_{2}}^2 - {s_{3}}^2 & s_{1}\, s_{2} & s_{1}\, s_{3}\\ s_{1}\, s_{2} &  - {s_{1}}^2 - {s_{3}}^2 & s_{2}\, s_{3}\\ s_{1}\, s_{3} & s_{2}\, s_{3} &  - {s_{1}}^2 - {s_{2}}^2 \end{array}\right]
\end{equation}
we get 
\begin{equation}
    \begin{aligned}
        \Tr(R) &= \Tr(\bm{I}) + \cancel{\Tr(\sin\theta(\bm{s}\cross))} + \Tr((1-\cos\theta)(\bm{s}\cross)^2)\\
        \implies r_{11} + r_{22} + r_{33} &= 3 +(1-\cos\theta)\underbrace{(-2s_1^2 -2s_2^2 -2s_3^2)}_{=-2}\\
        \implies \cos\theta &= \frac{r_{11} + r_{22} + r_{33} -1}{2}
    \end{aligned}
\end{equation}
and
\begin{equation}
    \begin{aligned}
        \begin{bmatrix}
            r_{11} & r_{12} & r_{13} \\  r_{21} & r_{22} & r_{23} \\  r_{31} & r_{32} & r_{33}
        \end{bmatrix} &= \begin{bmatrix}
            1 & 0& 0\\ 0&1 & 0 \\ 0& 0& 1
        \end{bmatrix} + 
        \sin\theta\begin{bmatrix}
            0 & -s_3 & s_2\\
        s_3 & 0 & -s_1\\
        -s_2 & s_1 & 0
        \end{bmatrix}\\ &+ 
        (1-\cos\theta) \begin{bmatrix}
             - {s_{2}}^2 - {s_{3}}^2 & s_{1}\, s_{2} & s_{1}\, s_{3}\\
             s_{1}\, s_{2} &  - {s_{1}}^2 - {s_{3}}^2 & s_{2}\, s_{3}\\
             s_{1}\, s_{3} & s_{2}\, s_{3} &  - {s_{1}}^2 - {s_{2}}^2
        \end{bmatrix} \\
        \implies &\begin{cases}
            r_{21} &= \sin\theta s_3 + (1-\cos\theta)s_1s_2\\
            r_{12} &= -\sin\theta s_3 + (1-\cos\theta)s_1s_2\\
            r_{13} &= \sin\theta s_2 + (1-\cos\theta)s_1s_3\\
            r_{31} &= -\sin\theta s_2 + (1-\cos\theta)s_1s_3\\
            r_{32} &= \sin\theta s_1 + (1-\cos\theta)s_2s_3\\
            r_{23} &= -\sin\theta s_1 + (1-\cos\theta)s_2s_3\\
        \end{cases}
    \end{aligned}
\end{equation}
By doing subtractions we get 
\begin{equation}
    \begin{aligned} 
        2\sin\theta\begin{bmatrix}
              s_1\\s_2\\s_3
        \end{bmatrix} &=
        \begin{bmatrix}
            r_{32} - r_{23}\\
            r_{13} - r_{31}\\
            r_{21} - r_{12}
        \end{bmatrix}
    \end{aligned}
\end{equation}
and connect this to three equations on the form
\begin{equation}
    \begin{aligned}
        \tan\theta = \frac{\sin\theta}{\cos\theta} &= \frac{\tfrac{r_{32} - r_{23}}{2s_1}}{\tfrac{r_{11} + r_{22} + r_{33} -1}{2}} = \frac{r_{32} - r_{23}}{s_1(r_{11} + r_{22} + r_{33} -1)}\\
        &= \frac{\tfrac{r_{13} - r_{31}}{2s_2}}{\tfrac{r_{11} + r_{22} + r_{33} -1}{2}} = \frac{r_{13} - r_{31}}{s_2(r_{11} + r_{22} + r_{33} -1)}\\
        &= \frac{\tfrac{r_{21} - r_{12}}{2s_3}}{\tfrac{r_{11} + r_{22} + r_{33} -1}{2}} = \frac{r_{21} - r_{12}}{s_3(r_{11} + r_{22} + r_{33} -1)}\\
        \implies \bm{s} &= \underline{\frac{1}{\tan\theta(r_{11} + r_{22} + r_{33} -1)} \begin{bmatrix}
            r_{32} - r_{23}\\
            r_{13} - r_{31}\\
            r_{21} - r_{12}
        \end{bmatrix}}
    \end{aligned}
\end{equation}

In the special case of $\sin\theta = 0$ we know that $\theta = n\pi$. For $n$ even we get $\me{\theta \bm{s}\cross}v = v$, and for $n$ odd we get that $\me{\theta \bm{s}\cross}v = -v$. This applies to all $\bm{s}$ and all $\theta=n\pi$ where $n \in \mathbb{Z}$.


\section{Trajectory generation}

Using the fuction \mcode{forward_kinematics} defined in \autoref{lst:2_1} we get the initial end effector position and orientation with 
$q(0) = \begin{bmatrix}
    0 & 0 & 90 & 0 &90 &0
\end{bmatrix}$ given as 
\begin{equation}
    C_6 = \begin{bmatrix}
        0 & 0 & 1 \\ 0 & -1 & 0 \\ 1 & 0 & 0
    \end{bmatrix} \quad
    \undertilde{o_6} = \begin{bmatrix}
        471.48 \\ 149.09 \\ 433.07
    \end{bmatrix}
\end{equation}

\lstinputlisting[style=Matlab-editor, caption=MATLAB code to do forward kinematics, captionpos = b, label=lst:2_1]{../../MATLAB/Assignment3/forward_kinematics.m}

Further we should use equation (25) and (26) from the notes, but we couldn't find a way to do this properly stepwise or without using inverse kinematics on wach step. We ended up with circular dependencies all the way, but have sort of a framework for the matlab script, given in \autoref{lst:2_2}. So this is how far we got.

\lstinputlisting[style=Matlab-editor, caption=MATLAB framwork for assignment 3, captionpos = b, label=lst:2_2]{../../MATLAB/Assignment3/assignment3.m}
\end{document}